% SPDX-License-Identifier: CC-BY-SA-4.0
% Part: Syntactic analysis
% Section: Regular grammmars
% Challenge exercise: Counter-automaton

\begingroup

\begin{exercice}[Mini-défi -- Automates à compteur]\label{exo:parsing/rational/challenge}

  Un automate à compteur est un automate à pile non déterministe dont l'alphabet de pile est limité à un seul
  symbole $X$  (en plus du symbole de fond de pile $\blank$ qui n'est jamais ajouté directement). Ainsi, la
  seule information qui peut être encodée dans la pile est sa taille (le compteur entier).
    
  Soit $\Sigma = \{a, b, c, d\}$. On considère la classe $\textsc{rec\_cpt}_\Sigma$ des langages
  sur $\Sigma$ reconnaissables par automate à compteur. Le but de cet exercice est de démontrer que
  $\textsc{rat}_\Sigma \subsetneq \textsc{rec\_cpt}_\Sigma \subsetneq \textsc{alg}_\Sigma$. 

  \begin{question}
  \item Justifier que $\textsc{rat}_\Sigma \subseteq \textsc{rec\_cpt}_\Sigma \subseteq \textsc{alg}_\Sigma$. 
  \item Montrer que $\left\{ a^n b^n \mid n \in \mathbb{N} \right\}$ est reconnaissable par automate à compteur.
  \item En déduire que $\textsc{rat}_\Sigma \neq \textsc{rec\_cpt}_\Sigma$. 
  \end{question}

  On veut maintenant démontrer que $\textsc{rec\_cpt}_\Sigma \neq \textsc{alg}_\Sigma$.
  Pour cela, on pose $L$ le langage engendré par la grammaire $\langle \Sigma, \{S\}, S, \{ S \rightarrow a S b \mid c S d \mid \epsilon \} \rangle$,
  et on veut montrer que $L \notin \textsc{rec\_cpt}_\Sigma$.

  Supposons (par contradiction) que $L \in \textsc{rec\_cpt}_\Sigma$, et posons $A$ un automate à compteur qui reconnaît $L$. 
  \begin{question}
  \item Montrer que pour tout mot $u \in L$, il existe un unique mot $v \in \mathcal{L}((a\mid c)^\star)$
    et un unique mot $\tilde{v} \in \mathcal{L}((b\mid d)^\star)$ tel que $u = v \cdot \tilde{v}$.
  \item Pour un $m \in \mathbb{N}$ donné, combien existe-t-il de mots $L$ de longueur $2m$ ? 
  \item Montrer qu'il existe $k \in \mathbb{N}$ tel que pour tout mot $u = v \cdot \tilde{v} \in L$,
    il existe une exécution de $A$ reconnaissant $u$ et passant par une configuration $\langle \tilde{v}, q, X^n\blank \rangle$
    telle que $n \le k |v|$. 
    
    Indice : On pourra poser, pour toute paire d'états $\langle q, q' \rangle$ et pour tout caractère $\sigma \in \Sigma$,
    $k_{q, a, q'}$ le nombre minimal de $X$ ajoutés à la pile dans un chemin menant de $q$ à $q'$, et suivant des $\varepsilon$-transitions puis
    une transition consommant $\sigma$, si un tel chemin existe. Ensuite, $k$ sera le plus grand des $k_{q, a, q'}$.
  \item Pour un $m$ donné, combien existe-t-il de configurations de la forme $\langle \varepsilon, q, X^n\blank \rangle$, telles que $n \le k m$ ?
  \item En déduire qu'il existe deux mots différents $v_1 \cdot \tilde{v_1} \neq v_2 \cdot \tilde{v_2} \in L$ tels que
    $v_1 \cdot \tilde{v_2} \in L$. 
  \item Conclure que $\textsc{rec\_cpt}_\Sigma \neq \textsc{alg}_\Sigma$.
  \end{question}
  
\end{exercice}

\begin{correction}

  \begin{question}

  \item
    \begin{itemize}
    \item $\textsc{rat}_\Sigma \subseteq \textsc{rec\_cpt}_\Sigma$ : tout automate fini se simule par un automate à compteur qui n'utilise jamais la pile (pile inchangée partout).
    \item $\textsc{rec\_cpt}_\Sigma \subseteq \textsc{alg}_\Sigma$ : tout automate à compteur est un automate à pile, donc reconnaît un langage algébrique.
    \end{itemize}

  \item Automate à compteur reconnaissant $\{a^n b^n \mid n\in\mathbb{N}\}$ :
    \begin{tikzpicture}[pushdown, baseline=(s.base)]
      \state[initial]   (s) at (0,0) {$0$};
      \state            (p) at (1,0) {$1$};
      \state[accepting] (f) at (2,0) {$2$};

      \path (s) edge[loop above] node {\smPAtrans{a}{\varepsilon}{X}}                      (s);
      \path (s) edge             node {\smPAtrans{\varepsilon}{\varepsilon}{\varepsilon}} (p);
      \path (p) edge[loop above] node {\smPAtrans{b}{X}{\varepsilon}}                      (p);
      \path (p) edge             node {\smPAtrans{\varepsilon}{\blank}{\varepsilon}}     (f);
    \end{tikzpicture}

  \item $\textsc{rat}_\Sigma \neq \textsc{rec\_cpt}_\Sigma$ (car $\{a^n b^n \mid n\in\mathbb{N}\}$ n'est pas rationnel).

  \end{question}

  \begin{question}

  \item En particulier, si $u\in L$ est engendré en appliquant $m$ fois une règle $aSb$ ou $cSd$,
    alors $v$ est le mot des $m$ lettres choisies dans $\{a,c\}$ (dans l'ordre), et $\tilde v$ est le mot des $m$ lettres correspondantes
    dans $\{b,d\}$ (dans l'ordre), avec la correspondance $a\leftrightarrow b$ et $c\leftrightarrow d$.

  \item Pour $m\in\mathbb{N}$, $|L\cap \Sigma^{2m}| = 2^m$.

  \item On fixe un automate à compteur $A=\langle \Sigma,\{X\},Q,I,F,\blank,\rightarrow\rangle$ reconnaissant $L$.
    Pour $q,q'\in Q$ et $\sigma\in\Sigma$, on définit $k_{q,\sigma,q'}$ comme le minimum, sur tous les chemins de transitions
    \[
      q \xrightarrow{\varepsilon/\ast/\ast}\cdots\xrightarrow{\varepsilon/\ast/\ast} q'' \xrightarrow{\sigma/\ast/\ast} q'
    \]
    (suite d'$\varepsilon$-transitions puis une transition consommant $\sigma$), du nombre total de symboles $X$ \emph{ajoutés} à la pile le long du chemin,
    si un tel chemin existe, et on pose $k_{q,\sigma,q'}=0$ sinon.
    On pose ensuite
    $k \eqdef \max\{k_{q,\sigma,q'} \mid q,q'\in Q,\ \sigma\in\Sigma\}$.
    Alors :
    $\forall u=v\cdot \tilde v\in L,\ \exists$ une exécution acceptante de $A$ passant par une configuration
      $\langle \tilde v,\ q,\ X^n\blank\rangle$ avec $n\le k|v|$.
    (On choisit une exécution acceptante et on considère l'instant où l'entrée restante devient $\tilde v$ ; le nombre de $X$ empilés pendant la lecture de $v$
    est majoré par $k$ par symbole lu, donc par $k|v|$.)

  \item Pour un $m$ donné, le nombre de configurations de la forme $\langle \varepsilon,\ q,\ X^n\blank\rangle$ telles que $n\le km$ est
    \[
      |Q|\,(km+1).
    \]

  \item Soit $m$ tel que $2^m > |Q|\,(km+1)$.
    Pour chaque $v\in (a\mid c)^m$, il existe un unique $\tilde v\in (b\mid d)^m$ tel que $v\tilde v\in L$.
    Donc il existe $2^m$ mots distincts $u=v\tilde v\in L\cap \Sigma^{2m}$.
    Par la question précédente et la question 3, pour chaque tel $u$ il existe une exécution acceptante atteignant une configuration
    $\langle \varepsilon,\ q,\ X^n\blank\rangle$ avec $n\le km$ après lecture complète de $u$.
    Par le principe des tiroirs, il existe deux mots distincts $u_1=v_1\tilde v_1\neq u_2=v_2\tilde v_2$ et deux exécutions acceptantes atteignant
    exactement la même configuration finale
    \[
      \langle \varepsilon,\ q,\ X^n\blank\rangle.
    \]
    Alors l'exécution de $u_1$ jusqu'à cette configuration, suivie de la \emph{fin} de l'exécution de $u_2$ (qui ne lit plus de symbole), donne une exécution
    acceptante pour $v_1\tilde v_2$.
    Donc
    \[
      v_1\tilde v_2 \in L.
    \]
    Or $u_1\neq u_2$ implique $v_1\neq v_2$, donc $\tilde v_1\neq \tilde v_2$ (unicité question 1), donc $v_1\tilde v_2\notin L$.
    Absurde.

  \item Donc $L\notin \textsc{rec\_cpt}_\Sigma$, donc $\textsc{rec\_cpt}_\Sigma \neq \textsc{alg}_\Sigma$.

  \end{question}

\end{correction}

\endgroup
\endinput
